% Concept map for the Conclusion: the whole chapter on one page.
% Replaces the "Key points: What this chapter has shown" box (15 Nov 2026).
% Hub = the exchange rate; six branches answer the chapter's questions,
% each tagged with the sections that cover it; the crisis box sits below.
\providecommand{\mapbox}[4][52mm]{% [width]{colour}{title}{body}
  \begin{tcolorbox}[enhanced, width=#1, arc=1.6mm, boxrule=0.9pt, colframe=#2, colback=white,
    colbacktitle=#2, coltitle=white, fonttitle=\sffamily\bfseries\footnotesize,
    fontupper=\sffamily\scriptsize\linespread{1.05}\selectfont, title={#3},
    left=4pt, right=4pt, top=3pt, bottom=3pt, toptitle=1.6pt, bottomtitle=1.6pt, boxsep=1pt]
  \RaggedRight #4\end{tcolorbox}}
\providecommand{\maptag}[1]{\hfill{\mdseries\scriptsize #1}}
\definecolor{mapgold}{HTML}{A9791A}
\definecolor{mapred}{HTML}{A23B2E}
\begin{tikzpicture}[
  font=\sffamily\scriptsize, line cap=round,
  box/.style={inner sep=0pt},
  link/.style={line width=1.3pt, #1},
]
% ---- hub ----
\node[fill=fxnavy, text=white, rounded corners=3mm, text width=34mm, align=center,
      inner sep=7pt, minimum height=27mm] (hub) at (0,0)
  {{\heavysans\large The exchange rate}\\[4pt]
   {\footnotesize ₹ per US dollar}\\[5pt]
   {\scriptsize the price of one currency\\ in terms of another}};

% ---- left column ----
\node[box] (L1) at (-5.75, 3.55) {\mapbox{fxteal}{Why it has a price\maptag{§1--2}}{%
  Goods, services and savings cross borders, so one currency must be swapped for another.
  India quotes rupees per dollar: a \textbf{bigger number} means a \textbf{weaker rupee}.}};
\node[box] (L2) at (-5.75, 0) {\mapbox{fxnavy}{What sets the price\maptag{§4--5, 8}}{%
  The demand for and supply of dollars meet at the \textbf{equilibrium rate}.\\[2pt]
  More rupees per dollar = \textbf{depreciation}. Fewer = \textbf{appreciation}.\\[2pt]
  The real rate and PPP adjust for prices.}};
\node[box] (L3) at (-5.75, -3.55) {\mapbox{fxplum}{How it is managed\maptag{§6--7}}{%
  Fixed, floating or \textbf{managed float}. India has used a managed float since 1993.\\[2pt]
  The RBI sells dollars from its \textbf{reserves} to slow a fall, and buys them to slow a rise.}};

% ---- right column ----
\node[box] (R1) at (5.75, 3.55) {\mapbox{fxorange}{Who buys and sells\maptag{§2--3}}{%
  \textbf{Demand:} importers, students abroad, Indians investing abroad.\\[2pt]
  \textbf{Supply:} exporters, remittances, foreign tourists and investors.\\[2pt]
  Every flow is recorded in the \textbf{balance of payments}.}};
\node[box] (R2) at (5.75, 0) {\mapbox{fxgreen}{What moves it\maptag{§9}}{%
  \textbullet\ Inflation above US inflation: a \textbf{slow drift}\\
  \textbullet\ US interest rates and capital flows: \textbf{sudden jumps}\\
  \textbullet\ Oil prices and the current account deficit\\
  \textbullet\ Expectations and speculation}};
\node[box] (R3) at (5.75, -3.55) {\mapbox{mapgold}{Who gains, who loses\maptag{§5, 10}}{%
  A weaker rupee \textbf{helps} exporters and families receiving remittances.\\[2pt]
  It \textbf{hurts} importers, students abroad and dollar borrowers, and pushes up prices
  (imported inflation).}};

% ---- crisis box below ----
\node[box] (C) at (0, -6.45) {\mapbox[115mm]{mapred}{When it breaks: currency crises\maptag{§11}}{%
  \textbf{Low reserves + heavy foreign debt + a large current account deficit} $\rightarrow$
  the currency collapses. India 1991 \textperiodcentered{} Sri Lanka 2022 \textperiodcentered{}
  Pakistan 2022--23. Today India's reserves cover about 11 months of imports.}};

% ---- links ----
\begin{scope}[on background layer]
  \draw[link=fxteal]   (hub.150) to[out=160, in=0] (L1.east);
  \draw[link=fxnavy]   (hub.west) -- (L2.east);
  \draw[link=fxplum]   (hub.210) to[out=200, in=0] (L3.east);
  \draw[link=fxorange] (hub.30)  to[out=20, in=180] (R1.west);
  \draw[link=fxgreen]  (hub.east) -- (R2.west);
  \draw[link=mapgold]  (hub.330) to[out=340, in=180] (R3.west);
  \draw[link=mapred]   (hub.south) -- (C.north);
\end{scope}

% ---- where next ----
\node[below=2.5mm of C, font=\sffamily\scriptsize, text=fxmuted, align=center]
  {\textbf{Where next?} The balance of payments, monetary policy and international trade all build on these ideas.};
\end{tikzpicture}
